
\documentclass[11pt,a4paper]{article}

\usepackage[a4paper,margin=0.9in]{geometry}
\usepackage[T1]{fontenc}
\usepackage{lmodern}
\usepackage{microtype}
\usepackage{setspace}
\usepackage{titlesec}
\usepackage{enumitem}
\usepackage{booktabs}
\usepackage{tabularx}
\usepackage{fancyhdr}
\usepackage{hyperref}

\hypersetup{hidelinks}
\setstretch{1.08}
\setlength{\parindent}{0pt}
\setlength{\parskip}{5pt}

\titleformat{\section}{\large\bfseries}{\thesection}{0.6em}{}
\titleformat{\subsection}{\normalsize\bfseries}{\thesubsection}{0.6em}{}

\pagestyle{fancy}
\fancyhf{}
\lhead{\small FinGuard AI}
\rhead{\small Prototype Document}
\cfoot{\small \thepage}
\renewcommand{\headrulewidth}{0.5pt}

\begin{document}

\begin{titlepage}
\centering
\vspace*{1.5cm}

{\Large\bfseries SOFTWARE ENGINEERING PROJECT\par}
\vspace{0.8cm}
{\Huge\bfseries FINGUARD AI\par}
\vspace{0.25cm}
{\Large Real-Time Financial Fraud Detection System\par}

\vspace{0.9cm}
\rule{0.75\textwidth}{0.6pt}\par
\vspace{0.4cm}
{\Large\bfseries PROTOTYPE DOCUMENT\par}
\vspace{0.4cm}
\rule{0.75\textwidth}{0.6pt}\par

\vfill

\begin{tabular}{rl}
\textbf{Prepared by:} & Arnav Kathuria (1024030068)\\
& Akshit Shandilya (1024030067)\\[0.2cm]
\textbf{Department:} & Computer Engineering\\
\textbf{Institution:} & Thapar Institute of Engineering and Technology\\[0.2cm]
\textbf{Date:} & September 2026
\end{tabular}

\vfill
\end{titlepage}

\section{Prototype Overview}

FinGuard AI is a real-time financial fraud detection system designed to evaluate transactions
as they occur and produce an explainable risk decision. The prototype combines deterministic
rules with supervised and unsupervised machine learning in a single scoring pipeline.

The basic prototype flow is:

\begin{center}
\textbf{Transaction}
$\rightarrow$
\textbf{Authentication}
$\rightarrow$
\textbf{Validation}
$\rightarrow$
\textbf{Feature Extraction}
$\rightarrow$
\textbf{Fraud Detection}
$\rightarrow$
\textbf{Risk Score}
$\rightarrow$
\textbf{Decision}
\end{center}

The system is intended to reduce detection delay, improve fraud detection quality, reduce
unnecessary false positives, and provide analysts with evidence explaining risky transactions.

\section{Objectives}

\begin{itemize}[leftmargin=1.4em,itemsep=2pt]
\item Evaluate transactions in real time.
\item Combine a rule engine, XGBoost, and Isolation Forest for hybrid fraud detection.
\item Produce a final risk score between 0 and 100.
\item Provide SHAP-based explanations for model-driven decisions.
\item Provide an analyst interface for reviewing flagged transactions.
\item Record analyst decisions for future model improvement.
\end{itemize}

\section{Prototype Workflow}

\subsection{Transaction Processing}

A customer submits a payment or transfer through the application. The API gateway authenticates
the request using JWT and applies role-based access control. The transaction service validates
and logs the transaction before extracting the features required for fraud detection.

The principal features include:

\begin{itemize}[leftmargin=1.4em,itemsep=2pt]
\item Transaction amount
\item Transaction velocity
\item Geographic deviation from customer history
\item Device fingerprint
\item Other available behavioural signals
\end{itemize}

\subsection{Fraud Detection}

The extracted features are evaluated by three detection components:

\begin{center}
\begin{tabularx}{0.92\textwidth}{>{\bfseries}p{0.27\textwidth}X}
\toprule
Component & Prototype role\\
\midrule
Rule Engine & Applies deterministic, hand-written fraud thresholds.\\
XGBoost & Produces a supervised probability-of-fraud signal.\\
Isolation Forest & Identifies unusual behavioural patterns and outliers.\\
\bottomrule
\end{tabularx}
\end{center}

The outputs are combined by a risk aggregation layer to generate a final score from 0 to 100.

\subsection{Risk Decision}

\begin{center}
\begin{tabular}{>{\bfseries}c c p{0.48\textwidth}}
\toprule
Risk tier & Score & Action\\
\midrule
Low & 0--30 & Approve transaction.\\
Medium & 31--60 & Hold and request customer verification.\\
High & 61--80 & Hold transaction and alert analyst.\\
Critical & 81--100 & Block transaction and open investigation case.\\
\bottomrule
\end{tabular}
\end{center}

\section{Prototype Interfaces}

\subsection{Customer Interface}

The customer interface provides a simple transaction form containing the required transaction
details. After submission, the system displays the resulting transaction status.

\begin{itemize}[leftmargin=1.4em,itemsep=2pt]
\item Transaction amount and recipient details
\item Submit transaction control
\item Transaction status
\item Verification message when additional confirmation is required
\end{itemize}

\subsection{Analyst Dashboard}

The analyst dashboard displays flagged transactions and allows an analyst to open an individual
case. A case contains:

\begin{itemize}[leftmargin=1.4em,itemsep=2pt]
\item Transaction identifier and amount
\item Final risk score and risk tier
\item Rule, XGBoost, and anomaly signals
\item Relevant transaction and behavioural features
\item SHAP feature contributions
\item Current case status
\item Analyst decision
\end{itemize}

The analyst can record one of the following outcomes:

\begin{center}
\textbf{Confirm Fraud}
\qquad
\textbf{False Positive}
\qquad
\textbf{Escalate Case}
\end{center}

\section{Explainability and Feedback}

SHAP is used to provide per-transaction feature attribution for model-driven decisions. Instead
of displaying only a risk score, the prototype presents the features that contributed to the
decision, such as transaction amount, device information, location deviation, velocity, and
account history.

Analyst decisions are stored in a decision feedback database. Confirmed fraud and false-positive
outcomes can subsequently support relabelling and future model retraining.

\begin{center}
\textbf{Detection}
$\rightarrow$
\textbf{Analyst Review}
$\rightarrow$
\textbf{Verdict}
$\rightarrow$
\textbf{Feedback Database}
$\rightarrow$
\textbf{Future Retraining}
\end{center}

\section{Technology Stack}

\begin{center}
\begin{tabularx}{0.92\textwidth}{>{\bfseries}p{0.30\textwidth}X}
\toprule
Layer & Proposed technology\\
\midrule
Authentication & JWT and role-based access control\\
Backend & REST API and transaction/scoring services\\
Supervised ML & XGBoost\\
Unsupervised ML & Isolation Forest with scikit-learn\\
Explainability & SHAP\\
Database & Relational or document database\\
Frontend & Web-based customer and analyst interfaces\\
\bottomrule
\end{tabularx}
\end{center}

\section{Evaluation}

The prototype will be evaluated using the project's defined KPIs:

\begin{itemize}[leftmargin=1.4em,itemsep=2pt]
\item \textbf{Fraud Recall:} proportion of actual fraud correctly flagged.
\item \textbf{False Positive Rate:} proportion of legitimate transactions incorrectly flagged.
\item \textbf{Mean Time to Detect:} latency between transaction occurrence and risk-score generation.
\item \textbf{Analyst Productivity:} change in average investigation time using the dashboard.
\item \textbf{Fraud Loss Avoided:} estimated monetary value of prevented fraud.
\end{itemize}

\section{Prototype Deliverables}

The prototype will provide:

\begin{enumerate}[leftmargin=1.5em,itemsep=2pt]
\item Transaction ingestion with JWT authentication and RBAC.
\item Rule engine and baseline XGBoost fraud scoring.
\item Risk aggregation with a 0--100 score.
\item Isolation Forest anomaly detection.
\item Basic analyst dashboard for flagged transactions.
\item SHAP-based explanations.
\item Analyst decision and feedback storage.
\item KPI reporting for system evaluation.
\end{enumerate}

\section{Conclusion}

The FinGuard AI prototype provides a concise implementation of the proposed fraud-detection
lifecycle. It combines deterministic rules, supervised and unsupervised machine learning,
risk aggregation, explainability, and analyst feedback into a single system. The prototype
therefore provides the foundation for implementing and evaluating the complete FinGuard AI
system.

\end{document}
